% !PRINTED 88
\weberpage{88}

Um ein System von nothwendigen, hinreichenden und von einander unabhängigen Bedingungen zu erhalten, fassen wir die Fragestellung etwas präciser und fragen nach den Bedingungen:

\gesperrt{dass aus einem System von $m$ linearen, homogenen Gleichungen mit $n$ Unbekannten $\nu$ von den Unbekannten durch $n-\nu$ willkürlich bleibende vollkommen bestimmt werden können.}

Auch diese Frage ist in \S.\,24 eigentlich schon beantwortet. Es muss unter den $\nu$-reihigen Determinanten eine von Null verschieden sein, während die $(\nu+1)$reihigen alle verschwinden. Es genügt aber schon, wenn es von einer kleineren Anzahl der $(\nu+1)$reihigen Determinanten feststeht, dass sie verschwinden.

Nehmen wir an, die Unbekannten $x_{\nu+1},x_{\nu+2},\ldots,x_n$ sollen willkürlich bleiben, $x_1,x_2,\ldots,x_\nu$ durch sie bestimmt sein, und nehmen die Determinante:
\begin{equation}
 A=\begin{vmatrix}
 a_1^{(1)}&a_2^{(1)}&\cdots&a_\nu^{(1)}\\
 a_1^{(2)}&a_2^{(2)}&\cdots&a_\nu^{(2)}\\
 \vdots&\vdots&&\vdots\\
 a_1^{(\nu)}&a_2^{(\nu)}&\cdots&a_\nu^{(\nu)}
 \end{vmatrix}
\tag{1}
\end{equation}
als von Null verschieden an. Wir berechnen die Unbekannten $x_1,x_2,\ldots,x_\nu$ nach \S.\,24, (9) und bilden die Summen (11), deren Verschwinden besagt, dass das gegebene Gleichungssystem wirklich befriedigt ist. Die Bedingungen dafür werden also:
\begin{equation}
\begin{vmatrix}
 a_1^{(1)}&a_2^{(1)}&\cdots&a_\nu^{(1)}&a_h^{(1)}\\
 a_1^{(2)}&a_2^{(2)}&\cdots&a_\nu^{(2)}&a_h^{(2)}\\
 \vdots&\vdots&&\vdots&\vdots\\
 a_1^{(\nu)}&a_2^{(\nu)}&\cdots&a_\nu^{(\nu)}&a_h^{(\nu)}\\
 a_1^{(k)}&a_2^{(k)}&\cdots&a_\nu^{(k)}&a_h^{(k)}
\end{vmatrix}=0,
\tag{2}
\end{equation}
oder anders geschrieben:
\begin{equation}
 Aa_h^{(k)}-\sum_{i=1}^{\nu}\sum_{\mu=1}^{\nu}
 a_h^{(i)}a_\mu^{(k)}A_\mu^{(i)}=0,
\tag{3}
\end{equation}
und diese Bedingungen genügen auch. Die Gleichung (2) oder (3) ist aber identisch befriedigt, wenn $h=1,2\ldots\nu$ oder $k=1,2\ldots\nu$ ist, und giebt also für diese Werthe keine Bedingung für die Coëfficienten. Solche Bedingungen ergeben sich nur für
%CONT>>
